\documentclass[10pt, twocolumn, landscape, a4paper]{article}

\usepackage[margin=1in]{geometry}
\usepackage[utf8]{inputenc}
\usepackage[spanish,es-noshorthands]{babel}
\usepackage{amssymb, amsmath, amsbsy}
\usepackage{mathpazo}
\usepackage{graphicx} % 
\usepackage{tikz}
\usetikzlibrary{arrows.meta}
% \usetikzlibrary{angles,quotes}
\usepackage[pdftex]{hyperref}
\hypersetup{pdftitle={Ecuación Punto|Pendiente de una Recta},pdfauthor={Luis Carrillo Gutiérrez}}

\renewcommand{\familydefault}{\sfdefault}
\pagestyle{empty}

\begin{document}
\subsection*{Geometría analítica}

\subsubsection*{Ecuación Punto-Pendiente de una Recta} 

Recta que pasa por $\mathrm{P}$ cuya pendiente es $m$. \\

\begin{tikzpicture}
	\begin{scope}[gray,<->]
	\draw (-2.5,0) -- (2.5,0);
	\draw (0,-1.25) -- (0,2.5);
	\end{scope}
	
	\draw[line width=1.3] (-4, .25) -- (4.5, 2.5);
	\draw (-1, 1.035) circle (.05);
	\draw (-1.25, .5) node {$\mathrm{P}\;(x_{_{1}}\,,\;y_{_{1}})$};
	%\draw (2.75, 2.035) circle (.05);
	%\draw (2.25, 2.5) node {$\mathrm{B}\;(x_{_{2}}\,,\;y_{_{2}})$};
	\draw (4, 1) node {\scalebox{1.4}{$y - y_{_{1}}   = m\,(x - x_{_{1}})$}};
\end{tikzpicture}
\end{document}